\documentclass[11pt]{article}
\usepackage[margin=1in]{geometry}
\usepackage{setspace}
\setstretch{1.15}

\begin{document}

\begin{center}
    \Large\textbf{Statement of Significance}
\end{center}
\vspace{1em}

Electronic Health Record (EHR) chart reviews are highly dynamic, requiring clinicians to frequently pivot between disjoint clinical topics within a single session. Standard search systems fail to adapt to these abrupt shifts, suffering from ``latent-context pollution'' where stale query history actively degrades the retrieval of new, critical information. 

This study introduces \textbf{Geodesic Diagnostic Trajectories (GDT)}, a novel differential geometry framework that models evolving clinician intent as a trajectory on a Riemannian manifold. By mathematically formalizing topic pivots using a geodesic shift ratio, GDT elegantly suppresses irrelevant historical context with provable interference bounds. Coupled with a predictive prefetch engine, this approach reduces information retrieval latency by nearly 27\% compared to standard baselines, translating to an estimated time savings of 4.1 minutes per clinician during a peak-acuity 4-hour shift. 

Ultimately, GDT provides a rigorous, scalable foundation for adaptive clinical decision support. By isolating the structural limitations of sparse lexical retrieval in complex clinical workflows, this work not only accelerates operational EHR navigation but also clearly defines the pathway for integrating modern dense semantic retrieval systems in healthcare settings.

\end{document}
